% Rosary Yoga -- The Manual
% Source booklet trim: 5.5 x 8.5" (US Half-Letter / Digest).
% This source produces SINGLE-PAGE PDF at trim size; the Makefile imposes
% it 2-up onto US Letter landscape via pdfpages.
%
% Page count target: ~28-32 pages. Must be a multiple of 4 for saddle-stitching.
% Currently: cover(4) + at-a-glance(2) + week(1) + decade(1) + pendant(2)
%            + dice(5) + bosses(7) + AF grid(2) + Sanctum(1) + Vigil(1) + colophon(1)
%            = 27 -> padded to 28 with a blank.
%
% Black ink only. ASCII-safe source. Lightly commented for the operator's edits.

\documentclass[10pt]{article}

% ---- page geometry: 5.5 x 8.5 inch trim, tight margins ---------------------
\usepackage[paperwidth=5.5in,paperheight=8.5in,
            top=0.55in,bottom=0.55in,left=0.55in,right=0.55in,
            headheight=14pt]{geometry}

\usepackage[T1]{fontenc}
\usepackage[utf8]{inputenc}
\usepackage{lmodern}
\usepackage{microtype}
\usepackage{csquotes}        % \enquote{} for curly quotes
\usepackage{tabularx}
\usepackage{booktabs}
\usepackage{array}
\usepackage{multicol}
\usepackage{tikz}            % foredge tab markers
\usepackage{fancyhdr}
\usepackage{enumitem}
\usepackage{xcolor}          % only used for shades-of-gray, no color ink

\pagestyle{empty}
\setlength{\parindent}{0pt}
\setlength{\parskip}{4pt}

% ---- foredge tab marker ----------------------------------------------------
% Prints a small black bar at a vertical position on the foredge (right edge of
% the recto, left edge of verso). Stack them across tabs so the closed book
% shows a printed index. Position is 0.0 (top) to 1.0 (bottom).
\newcommand{\foredgetab}[2]{%
  % #1 = position 0..1 from top of page (e.g. 0.15 = top tab)
  % #2 = label text (printed sideways inside the bar)
  \begin{tikzpicture}[remember picture,overlay]
    \node[anchor=east,inner sep=0pt] at ([yshift=-#1\paperheight,xshift=0pt]current page.north east) {%
      \begin{tikzpicture}
        \fill[black] (0,0) rectangle (0.30in,0.55in);
        \node[white,font=\bfseries\sffamily\scriptsize,rotate=-90] at (0.15in,0.275in) {#2};
      \end{tikzpicture}};
  \end{tikzpicture}%
}

% ---- section heading style (no chapters, no numbered sections) -------------
\newcommand{\manualtitle}[1]{%
  {\bfseries\sffamily\Large\MakeUppercase{#1}\par}%
  \vspace{2pt}\hrule height 0.5pt\vspace{8pt}}

\newcommand{\subhead}[1]{%
  \vspace{6pt}{\bfseries\sffamily\small\MakeUppercase{#1}\par}\vspace{2pt}}

% ============================================================================
\begin{document}

% ============== COVER (page 1) ==============================================
\thispagestyle{empty}
\vspace*{1.2in}
\begin{center}
  {\Huge\bfseries\sffamily ROSARY YOGA}\\[6pt]
  {\Large\sffamily THE MANUAL}\\[1.2in]
  \rule{2in}{0.5pt}\\[6pt]
  {\sffamily\small VERSION 1.0 \quad \textbar \quad YEAR ONE}\\[1.6in]
  {\sffamily\footnotesize The body learns the cue.\\ The mind follows the bead.\\ The paper records the year.}
\end{center}
\newpage

% ============== INSIDE FRONT (page 2): contents / how-to-use ================
\manualtitle{Contents}
\begin{description}[leftmargin=0pt,itemsep=2pt]
  \item[\sffamily PROTOCOL AT A GLANCE] one spread.
  \item[\sffamily THE WEEK]              modality by day-of-week.
  \item[\sffamily THE DECADE]            bead-phase load multipliers.
  \item[\sffamily THE PENDANT]           seven onboarding weeks.
  \item[\sffamily THE DICE TABLES]       one page per modality, plus modifiers.
  \item[\sffamily THE BOSSES]            one page per Year-One encounter.
  \item[\sffamily THE AF PT GRID]        scoring for the operator's bracket.
  \item[\sffamily THE SANCTUM]           the Sunday close-out.
  \item[\sffamily THE VIGIL]             the morning ritual (pending).
\end{description}
\vfill
\subhead{How to use this Manual}
Carry it until the protocol is memorized. Consult it thereafter. The tabbed
foredge lets the operator open to a section without reading the page numbers.
Outdated versions are visibly outdated; reprint when the protocol changes.
\newpage

% ============== PROTOCOL AT A GLANCE (pages 3-4) ============================
\foredgetab{0.12}{GLANCE}
\manualtitle{Protocol at a Glance}

\subhead{The three layers}
\begin{enumerate}[leftmargin=14pt,itemsep=1pt]
  \item \textbf{The bead} fixes the load (phase + multiplier).
  \item \textbf{The day} fixes the modality (Joyful / Sorrowful / Glorious / Luminous).
  \item \textbf{The dice} rolls the specific exercise within the modality.
\end{enumerate}

\subhead{The week in one line}
\textsf{\small MON mobility \textbar{} TUE endurance \textbar{} WED strength
\textbar{} THU skill \textbar{} FRI endurance \textbar{} SAT recovery \textbar{} SUN Sanctum.}

\subhead{The bead in one line}
\textsf{\small intro $\to$ intro $\to$ build $\to$ deload $\to$ build $\to$
peak $\to$ peak $\to$ peak $\to$ taper $\to$ boss.}

\subhead{The year in one line}
\textsf{\small pendant (7 wks) $\to$ 5 decades (10 wks each) $\to$ closing
(1 wk). Total $\approx$ 62 beads.}

\vspace{8pt}
\subhead{What the dice do not touch}
\begin{itemize}[leftmargin=14pt,itemsep=0pt]
  \item Modality. The day fixes that.
  \item Load. The bead fixes that.
  \item Boss schedule. The decade fixes that.
\end{itemize}

\subhead{Sundays}
No work. Close the bead. Plan the next.
\newpage

\foredgetab{0.12}{GLANCE}
\manualtitle{Protocol at a Glance \textnormal{\small / continued}}

\subhead{Day shape}
\begin{enumerate}[leftmargin=14pt,itemsep=1pt]
  \item Open the app (or the Logbook). Read RELAY's morning header.
  \item Roll the d6. Record the modifier (if any).
  \item Vigil. Then the day's work.
  \item Log scores. Close the day.
  \item Evening: \texttt{practice.md} -- the nightly rosary.
\end{enumerate}

\subhead{Boss week shape}
\begin{enumerate}[leftmargin=14pt,itemsep=1pt]
  \item Pre-fight: bead 9 taper, then bead 10 conserve.
  \item Fight day: protocol per Codex.
  \item Score is the score. Log it.
  \item Sanctum: RELAY writes the post-engagement entry.
\end{enumerate}

\vfill
\textit{\footnotesize The codec is in your pocket. If you need it on paper,
paper is allowed. -- RELAY}
\newpage

% ============== THE WEEK (page 5) ===========================================
\foredgetab{0.20}{WEEK}
\manualtitle{The Week}

\begin{tabularx}{\linewidth}{@{}l l l X@{}}
\toprule
\textbf{Day} & \textbf{Mystery} & \textbf{Modality} & \textbf{Notes} \\
\midrule
Mon & Joyful    & Mobility   & Vigil + extended flow. \\
Tue & Sorrowful & Endurance  & Run intervals. The grind. \\
Wed & Glorious  & Strength   & Heavier resistance. Lower reps. \\
Thu & Luminous  & Skill      & Technique. The ladder. \\
Fri & Sorrowful & Endurance  & Calisthenics grinds. The carrying. \\
Sat & Joyful    & Recovery   & Walk. Breath. Light. \\
Sun & --        & \textbf{Sanctum} & Rest. Logbook. Plan next bead. \\
\bottomrule
\end{tabularx}

\vspace{14pt}
\textit{\footnotesize Two endurance days, two mobility days, one strength,
one skill. The lopsidedness is real and correct for the room.}
\newpage

% ============== THE DECADE (page 6) =========================================
\foredgetab{0.28}{DECADE}
\manualtitle{The Decade}

\begin{tabularx}{\linewidth}{@{}c l l X@{}}
\toprule
\textbf{Bead} & \textbf{Phase} & \textbf{Load} & \textbf{Note} \\
\midrule
1  & Intro   & 0.50 & Patterning. Light. \\
2  & Intro   & 0.60 & Body remembers. \\
3  & Build   & 0.70 & Climbing. \\
4  & Deload  & 0.50 & The rosary breathes. \\
5  & Build   & 0.80 & Climbing again. \\
6  & Peak    & 1.00 & Full foundation load. \\
7  & Peak    & 1.15 & Highest sustained. \\
8  & Peak    & 1.30 & Highest of the decade. \\
9  & Taper   & 1.10 & Volume drops, intensity held. \\
10 & Boss    & N/A  & The fight is the fight. \\
\bottomrule
\end{tabularx}

\vspace{14pt}
\textit{\footnotesize The multiplier scales sets, reps, distance, or time --
whatever the modality denominates. See the dice tables for what scales.}
\newpage

% ============== THE PENDANT (pages 7-8) =====================================
\foredgetab{0.36}{PENDANT}
\manualtitle{The Pendant}

Seven weeks of onboarding. The body is being installed in the practice.

\vspace{6pt}
\begin{tabularx}{\linewidth}{@{}c l X@{}}
\toprule
\textbf{Bead} & \textbf{Name} & \textbf{What happens} \\
\midrule
0 & Crucifix         & Week zero. Read this. Set up the room. No work. \\
1 & First Hail Mary  & Vigil only, daily, 7 days. \\
2 & Second Hail Mary & Vigil + a daily walk. \\
3 & Third Hail Mary  & + 5 pushups, 5 situps, 5 wall-sits. \\
4 & OF\textsubscript{1} -- Cross & First trial. 10 PU, 10 SU, quarter-mile walk. \\
5 & Centerpiece -- Mary & Mary's codex transmitted this week. \\
6 & OF\textsubscript{2} -- Threshold & Half AF PT test. Year One begins next bead. \\
\bottomrule
\end{tabularx}

\vspace{12pt}
\textit{\footnotesize The two pendant trials are not bosses -- they are
baselines, scored against nothing but themselves.}
\newpage

\foredgetab{0.36}{PENDANT}
\manualtitle{The Pendant \textnormal{\small / scoring}}

\subhead{OF\textsubscript{1} -- The Cross}
\begin{tabularx}{\linewidth}{@{}l X@{}}
PU       & 10 reps, slow. Form over speed. \\
SU       & 10 reps. Anchor the feet. \\
Walk     & 1/4 mile, brisk pace. Not a run. \\
Record   & Numbers go to Codex. No grade. \\
\end{tabularx}

\vspace{10pt}
\subhead{OF\textsubscript{2} -- The Threshold}
\begin{tabularx}{\linewidth}{@{}l X@{}}
PU       & to fatigue (1 min cap). \\
SU       & to fatigue (1 min cap). \\
Plank    & to fatigue. \\
Run      & 1/2 mile. Sustainable pace. \\
Record   & Numbers go to Codex. This is the baseline that Decade 1's boss compares against. \\
\end{tabularx}

\vspace{14pt}
\textit{\footnotesize The Threshold is not a fight. It is a measurement.
Treat it as such.}
\newpage

% ============== DICE: JOYFUL (page 9) =======================================
\foredgetab{0.44}{DICE}
\manualtitle{Dice -- Joyful (Mobility / Recovery)}

\begin{tabularx}{\linewidth}{@{}c l X@{}}
\toprule
\textbf{d6} & \textbf{Variant} & \textbf{Description} \\
\midrule
1 & Long flow           & Vigil + 25 min slow flow. Hips, shoulders, spine. \\
2 & Yin hold            & Vigil + 5 yin poses, 3 min each. \\
3 & Walk + breath       & Vigil + 30 min walk, nasal breathing. \\
4 & Mobility ladder     & Vigil + joint-by-joint circuits (ankle to neck). \\
5 & Restorative flow    & Vigil + supine flow. Floor-only. \\
6 & Operator's choice   & Any of the above. Whatever the body asks for. \\
\bottomrule
\end{tabularx}

\vspace{10pt}
\subhead{Saturday rule}
On Joyful Saturday, halve any duration. Recovery is the goal.

\subhead{Boss-week rule}
On Joyful days in a boss week, default to variant 5 (restorative). Conserve.
\newpage

% ============== DICE: SORROWFUL (page 10) ===================================
\foredgetab{0.44}{DICE}
\manualtitle{Dice -- Sorrowful (Endurance)}

\begin{tabularx}{\linewidth}{@{}c l X@{}}
\toprule
\textbf{d6} & \textbf{Variant} & \textbf{Description} \\
\midrule
1 & 400m repeats        & 8 x 400m at threshold pace, 90s rest. \\
2 & 800m repeats        & 5 x 800m at tempo pace, 2 min rest. \\
3 & 1.5-mile time trial & Single effort. Sustainable. Negative split last quarter. \\
4 & Pushup ladder       & 1-2-3-...-N-...-3-2-1. Stop one short of failure. \\
5 & Carry circuit       & Farmer carry equivalent (no equipment: heavy bag, books). \\
6 & AMRAP 20            & 10 PU + 10 SU + 10 squats; rounds in 20 min. \\
\bottomrule
\end{tabularx}

\vspace{10pt}
\subhead{Bead scaling}
Distance and reps multiply by the bead's load multiplier (see Decade page).
On bead 8 (peak 1.30x), 8 x 400m becomes 10 x 400m.
\newpage

% ============== DICE: GLORIOUS (page 11) ====================================
\foredgetab{0.44}{DICE}
\manualtitle{Dice -- Glorious (Strength)}

\begin{tabularx}{\linewidth}{@{}c l X@{}}
\toprule
\textbf{d6} & \textbf{Variant} & \textbf{Description} \\
\midrule
1 & Pullup progression  & Dead hang $\to$ scapular pull $\to$ negative $\to$ rep. \\
2 & Pistol-squat ladder & Assisted $\to$ box $\to$ full. 5 sets per leg. \\
3 & Pushup variation    & Archer / decline / diamond / tempo. 5 hard sets. \\
4 & Plank complex       & Front + side + side + reverse. 4 rounds, 45s each. \\
5 & Wall-handstand work & Hold + shoulder taps + walks. 5 efforts. \\
6 & Operator's choice   & A target the operator wants to push. \\
\bottomrule
\end{tabularx}

\vspace{10pt}
\subhead{Quality rule}
Glorious is the only single-modality day of the week. If form breaks, stop
the set. Quality over load. Always.
\newpage

% ============== DICE: LUMINOUS (page 12) ====================================
\foredgetab{0.44}{DICE}
\manualtitle{Dice -- Luminous (Skill)}

\begin{tabularx}{\linewidth}{@{}c l X@{}}
\toprule
\textbf{d6} & \textbf{Variant} & \textbf{Description} \\
\midrule
1 & L-sit progression       & Tuck $\to$ one-leg $\to$ full. Tabata format. \\
2 & Handstand progression   & Wall $\to$ chest-to-wall $\to$ free balance. \\
3 & Pistol-squat technique  & Box height descents. Tempo 3-1-3. \\
4 & Pullup technique        & Hollow body. Scap control. Half-reps to full. \\
5 & Crow / arm balance      & Frog $\to$ crow $\to$ side crow. Hold for time. \\
6 & Free-form skill         & Pick any skill. Practice it for 30 min. \\
\bottomrule
\end{tabularx}

\vspace{10pt}
\subhead{Tempo rule}
Luminous is slow. If the work feels fast, slow it down. The body is learning.
\newpage

% ============== DICE: MODIFIERS (page 13) ===================================
\foredgetab{0.44}{DICE}
\manualtitle{Dice -- Modifiers}

After the day's modality is rolled, an optional second d6 may apply a modifier.

\begin{tabularx}{\linewidth}{@{}c l X@{}}
\toprule
\textbf{d6} & \textbf{Modifier} & \textbf{Effect} \\
\midrule
1 & None             & Run the day as written. \\
2 & Eccentric tempo  & 3-second lowering on every rep. Lower volume. \\
3 & Single-arm/leg   & Unilateral where possible. Half the volume per side. \\
4 & Paused reps      & 2-second pause at the hardest position. \\
5 & Density block    & Same work, set rest timer to 60s. \\
6 & INSTINCT         & Override. Listen to the body. Walk if needed. \\
\bottomrule
\end{tabularx}

\vspace{10pt}
\subhead{Haunt}
On a natural \enquote{6} followed by another \enquote{6} (probability $\approx$2\%), a
beaten boss surfaces as a haunt. The operator may accept the re-fight or
decline. Either is the protocol.
\newpage

% ============== BOSSES SUMMARY (pages 14-20) ================================
% Seven boss summary pages -- placeholders pending the boss-name draft.

\newcommand{\bosspage}[3]{%
  \foredgetab{0.52}{BOSSES}%
  \manualtitle{Boss \textnormal{--} #1}%
  \subhead{Anchor}\par #2 \par%
  \subhead{Fight protocol}\par #3 \par%
  \vfill%
  \textit{\footnotesize See Codex for first-contact transmission and fight log.}%
  \newpage}

\bosspage{BOSS\_PENDANT\_TRIAL (OF\textsubscript{2})}%
  {Pendant -- Threshold week.}%
  {PU to fatigue. SU to fatigue. Plank to fatigue. 1/2 mile run.
   Score = (PU + SU) + plank seconds. Baseline only.}

\bosspage{BOSS\_DECADE\_1 (e.g.\ The Sentinel)}%
  {Decade 1 -- The Joyful / Annunciation.}%
  {Full AF PT test. Pushups (1 min). Situps (1 min). Plank (max). 1.5-mile run.
   Score = AF composite. Baseline scouting -- the score is the score.}

\bosspage{BOSS\_DECADE\_2 (e.g.\ The Mirror)}%
  {Decade 2 -- The Sorrowful / Visitation.}%
  {Full PT test. Win = beat any one Sentinel category.
   Same components: PU, SU, plank, 1.5-mile run.}

\bosspage{BOSS\_DECADE\_3 (e.g.\ The Anvil)}%
  {Decade 3 -- The Glorious / Nativity.}%
  {No run. PU, pullups, SU, plank -- each to fatigue.
   Score = (PU + SU) + (Pullup-equiv x 2) + plank seconds.}

\bosspage{BOSS\_DECADE\_4 (e.g.\ The Pacer)}%
  {Decade 4 -- The Luminous / Presentation.}%
  {1.5-mile run only. Target = AF \enquote{Excellent} pace for bracket.
   Score = run time vs.\ threshold.}

\bosspage{BOSS\_DECADE\_5 (e.g.\ The Standard)}%
  {Decade 5 -- The Final / Coronation.}%
  {Full PT test, scored to \enquote{Excellent} across all four components.}

\bosspage{BOSS\_CLOSING\_TRIAL}%
  {Closing week.}%
  {Operator's choice. Re-fight the boss that hurt the most.
   Same protocol as that fight. Compare scores.}

% ============== AF PT GRID (pages 21-22) ===================================
\foredgetab{0.62}{AF PT}
\manualtitle{AF PT Scoring Grid}

\subhead{Operator's bracket}
Fill in once. Refer to thereafter. Update on birthday.

\begin{tabularx}{\linewidth}{@{}l X@{}}
Age      & \rule[-2pt]{1.5in}{0.4pt} \\
Sex      & \rule[-2pt]{1.5in}{0.4pt} \\
Bracket  & \rule[-2pt]{1.5in}{0.4pt} \\
Tested   & \rule[-2pt]{1.5in}{0.4pt} \\
\end{tabularx}

\vspace{14pt}
\subhead{Thresholds (this bracket)}

\begin{tabularx}{\linewidth}{@{}l c c c c@{}}
\toprule
\textbf{Component} & \textbf{Min} & \textbf{Pass} & \textbf{Good} & \textbf{Excellent} \\
\midrule
Pushups (1 min)        & \_\_\_ & \_\_\_ & \_\_\_ & \_\_\_ \\
Situps (1 min)         & \_\_\_ & \_\_\_ & \_\_\_ & \_\_\_ \\
Plank (sec)            & \_\_\_ & \_\_\_ & \_\_\_ & \_\_\_ \\
Run 1.5 mi (mm:ss)     & \_\_\_ & \_\_\_ & \_\_\_ & \_\_\_ \\
Composite (out of 100) & \_\_\_ & 75    & 85    & 90    \\
\bottomrule
\end{tabularx}

\vspace{8pt}
\textit{\footnotesize Look up the operator's bracket in the current AF DAFMAN
36-2905 charts. Pencil in the numbers above. The Manual is then complete for
this year.}
\newpage

\foredgetab{0.62}{AF PT}
\manualtitle{AF PT Scoring Grid \textnormal{\small / conversions}}

\subhead{Pullup equivalence (Codex bosses)}
\begin{tabularx}{\linewidth}{@{}c X@{}}
\toprule
\textbf{Pullups} & \textbf{Equivalent} \\
\midrule
0--3   & Assisted: 3 negatives = 1 equiv. \\
4--10  & Direct count. \\
11+    & Direct count, weighted reps may count as 1.5x. \\
\bottomrule
\end{tabularx}

\vspace{10pt}
\subhead{Plank score band}
\begin{tabularx}{\linewidth}{@{}l c@{}}
\toprule
\textbf{Hold} & \textbf{Band} \\
\midrule
< 1:00         & Min \\
1:00 -- 2:00   & Pass \\
2:00 -- 3:30   & Good \\
3:30+          & Excellent \\
\bottomrule
\end{tabularx}

\vspace{10pt}
\textit{\footnotesize These bands are an internal aid -- they do not replace
the official AF chart. Use them when the chart isn't to hand.}
\newpage

% ============== SANCTUM (page 23) ==========================================
\foredgetab{0.74}{SANCTUM}
\manualtitle{The Sanctum}

\subhead{What it is}
The Sunday close-out. The most ritualized moment in the practice outside of
boss encounters. The operator does no work.

\subhead{The four moves}
\begin{enumerate}[leftmargin=14pt,itemsep=2pt]
  \item \textbf{The week in numbers.} Days completed, dice rolled, scores
    against target, weight, sleep.
  \item \textbf{The week in voice.} RELAY's post-bead transmission.
    Generated in app; transcribed by hand in the paper Logbook if used.
  \item \textbf{The Sanctum note.} A single text field. A sentence, a
    paragraph, nothing. The operator's hand in the record.
  \item \textbf{Close the bead.} Page-turn in the Logbook; primary button
    in the app. Commits the bead. Opens the next.
\end{enumerate}

\subhead{The Sanctum rule}
Sundays are not for work. Not catch-up. Not extra. If a day was missed,
RELAY logs it and the week moves on. The bead does not advance until Sunday.

\vfill
\textit{\footnotesize The Sanctum is the only place the operator is
encouraged to linger.}
\newpage

% ============== VIGIL (page 24) ============================================
\foredgetab{0.86}{VIGIL}
\manualtitle{The Vigil}

\subhead{What it is}
The morning warmup. Approximately ten minutes. Carpeted floor, no equipment,
no impact.

\subhead{What it must do}
\begin{itemize}[leftmargin=14pt,itemsep=1pt]
  \item Prepare the body for the day's modality.
  \item Be memorizable in a week. Run without a sheet.
  \item Leave the body warm and present, not drained.
  \item Reset spine, hips, shoulders.
\end{itemize}

\vspace{8pt}
\subhead{Specification \emph{pending}}
Three candidates are under review:
\begin{enumerate}[leftmargin=14pt,itemsep=1pt]
  \item Sun Salutation A (five rounds, step-back variant).
  \item Joe DeFranco's Agile 8.
  \item Original Strength RESETs.
\end{enumerate}

The recommendation will be folded in here once the comparison returns. Until
then, the operator may use any one of the three; consistency matters more
than which.

\vfill
\textit{\footnotesize \enquote{The hour is yours. Begin.} -- VIGIL}
\newpage

% ============== COLOPHON (page 25) =========================================
\manualtitle{Colophon}

This Manual is typeset in LaTeX with Latin Modern at 10pt. It is printed on
US Letter, duplex, two pages to a side, then folded in half and stapled
through the spine with a long-arm stapler. Finished trim: 5.5 x 8.5".

Source lives at \texttt{paper/manual/manual.tex}. Edit, recompile, reprint.
The Manual is meant to be reprinted whenever the protocol changes; outdated
versions become visibly outdated by their version number on the cover.

\vspace{16pt}
\subhead{Versions}
\begin{tabularx}{\linewidth}{@{}l l X@{}}
\toprule
\textbf{Ver.} & \textbf{Date} & \textbf{Change} \\
\midrule
1.0 & --- & Initial release. \\
    &    &  \\
    &    &  \\
\bottomrule
\end{tabularx}

\vfill
\textit{\footnotesize The body learns the cue. The mind follows the bead.
The paper records the year.}
\newpage

% ============== PADDING TO MULTIPLE OF 4 ====================================
% Currently at 25 content pages -- pad to 28 with three blank back-matter pages.
\thispagestyle{empty}\null\vfill
\begin{center}\textit{\footnotesize Notes}\end{center}
\vfill\newpage
\thispagestyle{empty}\null\vfill
\begin{center}\textit{\footnotesize Notes}\end{center}
\vfill\newpage

% ============== BACK COVER (page 28) =======================================
\thispagestyle{empty}
\vspace*{\fill}
\begin{center}
  \rule{2in}{0.5pt}\\[8pt]
  {\sffamily\small ROSARY YOGA}\\[2pt]
  {\sffamily\footnotesize THE MANUAL \quad v1.0}\\[1in]
  {\sffamily\scriptsize \enquote{The codec is in your pocket.\\
  If you need it on paper, paper is allowed.}}\\[6pt]
  {\sffamily\scriptsize -- RELAY}
\end{center}
\vspace*{\fill}

\end{document}
